\chapter{Ring dan Aritmetika}\label{ch:b3:rings}

Bilangan bulat biasa terfaktorkan secara tunggal atas bilangan prima; demikian pula
polinomial atas sebuah lapangan. Apakah kedua kenyataan itu satu teorema yang sama? Bab ini
menjawab ya, dan menemukan hipotesis persis yang membuat sebuah ``aritmetika''
mungkin di dalam ring komutatif, yakni rantai
\[
\text{Euclid} \;\Longrightarrow\; \text{ideal utama}
\;\Longrightarrow\; \text{faktorial (DFT)},
\]
dengan semua implikasinya dibuktikan dan semua konversnya disanggah. Teori ini
lalu diuji di tempat ia benar-benar berguna: bilangan bulat Gauss
$\Z[\iu]$ (yang akan memecahkan teorema dua kuadrat Fermat pada
soal akhir pekan), ring polinomial dalam beberapa variabel (lema
Gauss, kriteria \omterm{thm:b3:rings:criteria}{Eisenstein}), dan \omterm{def:b3:rings:noetherian}{ring Noether}, yang memuncak
pada teorema basis Hilbert. Sepanjang bab ini, \emph{ring} berarti
ring komutatif dengan unsur satuan $1 \neq 0$; \omterm{def:b3:rings:ideal}{ideal} $\Z$ dan $K[X]$
dari jilid Tahun ke-2 menjadi dua contoh penuntun kita.

\section{Ideal, kuosien, dan teorema isomorfisma}

\begin{definition}\label{def:b3:rings:ideal}
Sebuah \emph{ideal}\index{ideal} $I$ dari ring $A$ adalah subgrup
aditif yang memenuhi $AI \subseteq I$. Yang disebut \emph{ring
kuosien}\index{ring kuosien} $A/I$ adalah \omterm{thm:b3:groups:quotient}{grup kuosien} $(A, +)/I$
dengan perkalian $(a + I)(b + I) = ab + I$: perkalian ini terdefinisi dengan baik,
sebab mengubah $a$ menjadi $a + x$ dan $b$ menjadi $b + y$ ($x, y \in I$)
mengubah $ab$ sebesar $ay + xb + xy \in I$. Proyeksi $\pi \colon A
\to A/I$ adalah morfisma ring surjektif dengan kernel $I$, dan kernel
morfisma ring tidak lain adalah ideal.
\end{definition}

\begin{theorem}[Teorema isomorfisma pertama]\label{thm:b3:rings:firstiso}
Jika $f \colon A \to B$ morfisma ring, maka $\bar f\colon
A/\ker f \to \operatorname{im} f$, $a + \ker f \mapsto f(a)$, adalah
isomorfisma ring. Secara lebih umum, $f$ terfaktorkan melalui $A/I$ untuk sembarang
\omterm{def:b3:rings:ideal}{ideal} $I \subseteq \ker f$. \omterm{def:b3:rings:ideal}{Ideal} dari $A/I$ berupa $J/I$ dengan
$J \supseteq I$ sebuah \omterm{def:b3:rings:ideal}{ideal} $A$ (teorema korespondensi).
\end{theorem}

\begin{proof}
Sama seperti pada grup (\cref{thm:b3:groups:firstiso,%
thm:b3:groups:correspondence}), dengan mencatat bahwa semua pemetaan yang terlibat juga
menghormati perkalian: $\bar f$ terdefinisi dengan baik, bijektif pada
petanya, dan multiplikatif; korespondensi $J \mapsto J/I$,
$\bar J \mapsto \pi^{-1}(\bar J)$ mempertahankan \omterm{def:b3:rings:ideal}{ideal} dalam kedua
arah karena $\pi$ morfisma ring yang surjektif.
\end{proof}

\begin{definition}\label{def:b3:rings:primemaximal}
Misalkan $I \subsetneq A$ sebuah \omterm{def:b3:rings:ideal}{ideal} sejati. \omterm{def:b3:rings:ideal}{Ideal} $I$ disebut
\emph{prima}\index{ideal prima} jika $ab \in I \Rightarrow a \in I$
atau $b \in I$; dan $I$ disebut \emph{maksimal}\index{ideal maksimal} jika tak ada
\omterm{def:b3:rings:ideal}{ideal} yang terletak tepat di antara $I$ dan $A$.
\end{definition}

\begin{proposition}\label{prop:b3:rings:primemaximal}
$I$ prima $\iff$ $A/I$ daerah integral; $I$ \omterm{def:b3:rings:primemaximal}{maksimal}
$\iff$ $A/I$ lapangan. Khususnya, \omterm{def:b3:rings:primemaximal}{ideal maksimal} selalu prima.
\end{proposition}

\begin{proof}
Tulis $\bar a$ untuk kelas di $A/I$. Kalimat ``$I$ prima'' beralih
kata demi kata menjadi ``$\bar a\bar b = 0 \Rightarrow \bar a = 0$ atau $\bar b
= 0$'', dan $A/I \neq 0$ menjadi $I \neq A$: itulah definisi
daerah integral. Untuk kemaksimalan, pakailah teorema korespondensi: tidak ada \omterm{def:b3:rings:ideal}{ideal}
yang terletak tepat di antara $I$ dan $A$ $\iff$ $A/I$ tak punya \omterm{def:b3:rings:ideal}{ideal} selain
$0$ dan dirinya sendiri $\iff$ $A/I$ lapangan --- untuk langkah terakhir: di dalam
lapangan satu-satunya \omterm{def:b3:rings:ideal}{ideal} adalah $0$ dan seluruhnya (\omterm{def:b3:rings:ideal}{ideal} yang memuat
$x \ne 0$ memuat $x^{-1}x = 1$); sebaliknya jika setiap $x$ tak nol
membangun \omterm{def:b3:rings:ideal}{ideal} satuan, maka $xy = 1$ untuk suatu $y$. Lapangan adalah
daerah integral, jadi \omterm{def:b3:rings:primemaximal}{ideal maksimal} selalu prima.
\end{proof}

\begin{example}\label{ex:b3:rings:primemaximal}
Di $\Z$: \omterm{def:b3:rings:primemaximal}{ideal primanya} adalah $(0)$ dan semua $(p)$ dengan $p$ prima;
yang \omterm{def:b3:rings:primemaximal}{maksimal} adalah $(p)$ (sebab $\Z/p\Z = \mathbb F_p$ lapangan, sedangkan
$\Z/(0) = \Z$ bukan). Di $K[X, Y]$: $(X) \subsetneq (X, Y)$ keduanya
prima ($K[X,Y]/(X) \cong K[Y]$ daerah integral;
$K[X,Y]/(X,Y) \cong K$ lapangan), jadi $(X)$ prima tetapi tidak
\omterm{def:b3:rings:primemaximal}{maksimal}.
\end{example}

Untuk menjamin bahwa \omterm{def:b3:rings:primemaximal}{ideal maksimal} \emph{ada} dalam keumuman penuh,
kita memerlukan sebuah asas teori himpunan. Sebuah himpunan terurut parsial disebut
\emph{induktif} jika setiap himpunan bagian yang terurut total (sebuah \emph{rantai})
mempunyai batas atas.

\begin{theorem}[Lema Zorn]\label{thm:b3:rings:zorn}
Setiap himpunan terurut parsial yang induktif dan tak kosong mempunyai unsur
\omterm{def:b3:rings:primemaximal}{maksimal}.\index{lema Zorn}
\admitted
\end{theorem}

\begin{remark}\label{rem:b3:rings:zorn}
Pernyataan ini bukan teorema matematika biasa melainkan sebuah \emph{aksioma}:
di atas aksioma dasar Zermelo--Fraenkel dalam teori
himpunan, ia setara dengan aksioma pilihan (``setiap hasil kali himpunan tak kosong
juga tak kosong''), yang kita terima sepanjang buku ini. Setiap
pemakaiannya kita tandai. Analisis akan memanggilnya lagi (Hahn--Banach,
\cref{ch:b3:banach}).
\end{remark}

\begin{theorem}[Krull]\label{thm:b3:rings:krull}
Setiap \omterm{def:b3:rings:ideal}{ideal} sejati $I \subsetneq A$ termuat di dalam sebuah \omterm{def:b3:rings:primemaximal}{ideal
maksimal}.\index{teorema Krull}
\end{theorem}

\begin{proof}
Urutkan menurut inklusi himpunan $\mathcal E$ yang beranggotakan \omterm{def:b3:rings:ideal}{ideal} sejati
pemuat $I$; himpunan ini tak kosong ($I \in \mathcal E$). Sebuah rantai
$(J_\lambda)$ di $\mathcal E$ mempunyai batas atas $J = \bigcup
J_\lambda$: ia \omterm{def:b3:rings:ideal}{ideal} (dua unsur $a, b \in J$ terletak pada satu $J_\lambda$
yang sama menurut ketotalan), sejati ($1 \notin J_\lambda$ untuk setiap $\lambda$),
dan memuat $I$. Lema Zorn lalu memberikan unsur \omterm{def:b3:rings:primemaximal}{maksimal} $\mathcal
E$, yaitu \omterm{def:b3:rings:primemaximal}{ideal maksimal} pemuat $I$ (\omterm{def:b3:rings:ideal}{ideal} sejati yang terletak tepat
di atasnya akan menjadi anggota $\mathcal E$).
\end{proof}

\begin{theorem}[Teorema sisa Tiongkok]\label{thm:b3:rings:crt}
Misalkan $I_1, \dots, I_n$ \omterm{def:b3:rings:ideal}{ideal} $A$ yang \emph{komaksimal}\index{ideal
komaksimal} berpasangan ($I_k + I_l = A$ untuk $k \neq l$). Maka
\[
A\Big/\bigcap_{k=1}^n I_k \;\xrightarrow{\;\sim\;}\;
\prod_{k=1}^n A/I_k,
\qquad
a \longmapsto (a + I_1, \dots, a + I_n),
\]
dan lebih lanjut $\bigcap_k I_k = I_1 I_2 \cdots I_n$ (\omterm{def:b3:rings:ideal}{ideal} yang
dibangun oleh hasil kali).\index{teorema sisa Tiongkok}
\end{theorem}

\begin{proof}
Pemetaan $f(a) = (a + I_k)_k$ adalah morfisma ring dengan kernel
$\bigcap I_k$; menurut \cref{thm:b3:rings:firstiso} cukup
dibuktikan kesurjektifannya. Tetapkan $k$; untuk setiap $l \neq k$ tulis $1 = u_l +
v_l$ dengan $u_l \in I_k$, $v_l \in I_l$ (komaksimalitas). Maka
\[
e_k = \prod_{l \neq k} v_l = \prod_{l\neq k}(1 - u_l)
\equiv 1 \pmod{I_k},
\qquad e_k \in I_l \ (l \neq k),
\]
sehingga $f(e_k) = (0, \dots, 1, \dots, 0)$; diberikan sasaran $(a_k +
I_k)_k$, unsur $\sum_k a_k e_k$ terpetakan padanya.

Hasil kali lawan irisan: selalu berlaku $I_1\cdots I_n \subseteq \bigcap I_k$.
Sebaliknya, dengan induksi cukup ditangani kasus $n = 2$
(mudah diperiksa bahwa $I_1$ dan $I_2\cdots I_n$ \omterm{thm:b3:rings:crt}{komaksimal}: mengalikan
$1 = u_l + v_l$ atas $l \geq 2$ memberikan $1 \in I_1 + I_2\cdots
I_n$). Untuk $n = 2$: tulis $1 = u + v$, $u \in I_1$, $v \in I_2$;
untuk $x \in I_1 \cap I_2$ berlaku $x = xu + xv \in I_2I_1 + I_1I_2 =
I_1I_2$.
\end{proof}

\begin{example}\label{ex:b3:rings:crtz}
Di $\Z$ dengan $I_k = (m_k)$ dan $m_k$ relatif prima berpasangan: $\Z/(m_1\cdots
m_n)\Z \cong \prod \Z/m_k\Z$ --- inilah teorema sisa Tiongkok pada
jilid Tahun ke-2. Dibatasi pada unitnya: $(\Z/mn\Z)^\times \cong
(\Z/m\Z)^\times \times (\Z/n\Z)^\times$ untuk $\gcd(m,n)=1$, dan dari sini
diperoleh kemultiplikatifan fungsi $\varphi$ Euler
(\cref{exo:b3:rings:8}).
\end{example}

\section{Keterbagian: Euclid, ideal utama, faktorial}

\begin{definition}\label{def:b3:rings:divisibility}
Misalkan $A$ daerah integral dan $a, b \in A$. Kita katakan bahwa $a$
\emph{membagi} $b$ ($a \mid b$) jika $b \in (a) = aA$. Unsur $a, b$
disebut \emph{sekawan}\index{unsur sekawan} jika $a = ub$
dengan $u \in A^\times$ (ekuivalen dengan $(a) = (b)$). Unsur $p$ yang
tak nol dan bukan unit disebut:
\begin{itemize}
  \item \emph{tak tereduksi}\index{unsur tak tereduksi} jika $p = ab$
        memaksa $a \in A^\times$ atau $b \in A^\times$;
  \item \emph{prima}\index{unsur prima} jika $p \mid ab$ memaksa $p
        \mid a$ atau $p \mid b$ (yakni \omterm{def:b3:rings:ideal}{ideal} $(p)$ prima).
\end{itemize}
\end{definition}

\begin{proposition}\label{prop:b3:rings:primeirred}
Di dalam sembarang daerah integral, prima $\Rightarrow$ \omterm{def:b3:rings:divisibility}{tak tereduksi}. Kebalikannya
salah secara umum: di $\Z[\iu\sqrt 5] = \{a + \iu b\sqrt5 : a, b
\in \Z\}$, unsur $2$ \omterm{def:b3:rings:divisibility}{tak tereduksi} tetapi tidak prima.
\end{proposition}

\begin{proof}
Misalkan $p$ prima dan $p = ab$. Maka $p \mid ab$, katakanlah $p \mid a$:
$a = pc$, sehingga $p = pcb$, dan mencoret $p$ (di daerah integral!) memberikan $cb =
1$: jadi $b \in A^\times$.

Di $\Z[\iu\sqrt5]$, pakailah norma $N(x + \iu y\sqrt 5) = x^2 +
5y^2$, yang multiplikatif (sebab ia $\abs z^2$). Jika $2 = ab$
dengan $a, b$ bukan unit, maka $4 = N(a)N(b)$ dengan $N(a), N(b) \neq
1$ (unsur bernorma $1$ hanyalah $\pm1$, yaitu unitnya), sehingga $N(a) = 2$:
mustahil, sebab $x^2 + 5y^2 = 2$ tak punya penyelesaian bulat. Jadi $2$
\omterm{def:b3:rings:divisibility}{tak tereduksi}. Namun $2 \mid 6 = (1 + \iu\sqrt5)(1 - \iu\sqrt5)$
sedangkan $2$ tidak membagi satu pun faktornya (sebab $\frac12 \pm \frac{\iu\sqrt5}2
\notin \Z[\iu\sqrt5]$): jadi tidak prima.
\end{proof}

\begin{definition}\label{def:b3:rings:pidufd}
Sebuah daerah integral $A$ disebut:
\begin{itemize}
  \item \emph{Euclid}\index{daerah Euclid} jika ada pemetaan
        $\nu \colon A \setminus \{0\} \to \N$ (sebuah \emph{fungsi
        Euclid}) sedemikian sehingga untuk setiap $a, b$ dengan $b \ne 0$
        terdapat $q, r$ dengan $a = bq + r$ dan ($r = 0$ atau $\nu(r) <
        \nu(b)$);
  \item \emph{\omterm{def:b3:rings:ideal}{ideal} utama} (sebuah \emph{DIU}\index{daerah ideal
        utama}) jika setiap \omterm{def:b3:rings:ideal}{idealnya} berbentuk $(a)$;
  \item \emph{faktorial} (sebuah \emph{DFT}\index{daerah faktorisasi
        tunggal}) jika setiap unsur tak nol yang bukan unit merupakan hasil kali
        \omterm{def:b3:rings:divisibility}{unsur tak tereduksi}, secara tunggal sampai pada urutan dan kesekawanan.
\end{itemize}
\end{definition}

\begin{theorem}\label{thm:b3:rings:euclideanpid}
\omterm{def:b3:rings:pidufd}{Euclid} $\Rightarrow$ \omterm{def:b3:rings:ideal}{ideal} utama.
\end{theorem}

\begin{proof}
Misalkan $I \neq (0)$ sebuah \omterm{def:b3:rings:ideal}{ideal} dan $b \in I \setminus\{0\}$ dengan
$\nu(b)$ minimal. Untuk $a \in I$, lakukan pembagian: $a = bq + r$; maka $r = a
- bq \in I$, dan $\nu(r) < \nu(b)$ akan bertentangan dengan keminimalan, sehingga
$r = 0$ dan $a \in (b)$: jadi $I = (b)$.
\end{proof}

\begin{example}\label{ex:b3:rings:euclidean}
$\Z$ (dengan $\nu = \abs\cdot$) dan $K[X]$ (dengan $\nu = \deg$) bersifat
\omterm{def:b3:rings:pidufd}{Euclid} --- jilid Tahun ke-2 membuktikan kedua pembagian itu. Demikian pula
$\Z[\iu]$, dengan $\nu = N$ berupa norma kuadrat
(\cref{exo:b3:rings:4}); geometri pembuktiannya tampak pada
gambar di bawah. \omterm{def:b3:rings:pidufd}{DIU} yang bukan \omterm{def:b3:rings:pidufd}{Euclid} memang ada tetapi sulit
dipastikan (contoh bakunya adalah
$\Z\bigl[\frac{1+\iu\sqrt{19}}2\bigr]$); sedangkan \omterm{def:b3:rings:pidufd}{DFT} yang bukan \omterm{def:b3:rings:pidufd}{DIU}
mudah dicari: $K[X, Y]$ (\cref{exo:b3:rings:6}), atau $\Z[X]$.
\end{example}

\begin{omfigure}
\begin{tikzpicture}[scale=1.05]
  \clip (-2.7,-1.7) rectangle (2.7,2.2);
  \foreach \x in {-2,...,2}
    \foreach \y in {-1,...,2}
      \fill[black!55] (\x,\y) circle (1.4pt);
  \draw[omDef, thick] (0.6,0.4) circle (1);
  \fill[omThm] (0.6,0.4) circle (1.8pt)
    node[below right, font=\small, omThm] {$a/b$};
  \fill[omDef] (1,0) circle (2pt)
    node[below, font=\small, omDef] {$q$};
  \draw[omDef, ->, shorten >=2.5pt, shorten <=2.5pt]
    (0.6,0.4) -- (1,0);
\end{tikzpicture}

{\small Pembagian di dalam bilangan bulat Gauss: hasil bagi eksak
$a/b \in \C$ terletak pada jarak $\leq \frac{\sqrt2}{2} < 1$ dari
suatu titik kisi $q \in \Z[\iu]$; lalu $r = a - bq$ mempunyai $N(r) =
N(b)\,\abs{a/b - q}^2 < N(b)$. Satu pembagian \omterm{def:b3:rings:pidufd}{Euclid}, dan lahirlah
seluruh aritmetika.}
\end{omfigure}

\begin{lemma}[Rantai naik ideal utama]
\label{lem:b3:rings:acc}
Di dalam \omterm{def:b3:rings:pidufd}{DIU}, setiap barisan naik \omterm{def:b3:rings:ideal}{ideal} $I_1 \subseteq I_2
\subseteq \cdots$ akhirnya menjadi konstan.
\end{lemma}

\begin{proof}
$I = \bigcup_n I_n$ adalah \omterm{def:b3:rings:ideal}{ideal} (sebab gabungannya naik), jadi $I =
(a)$; unsur $a$ terletak pada suatu $I_N$, dan lalu $I = (a)
\subseteq I_N \subseteq I_n \subseteq I$ untuk $n \geq N$.
\end{proof}

\begin{lemma}[Bézout; lema Euclid]\label{lem:b3:rings:bezout}
Misalkan $A$ sebuah \omterm{def:b3:rings:pidufd}{DIU} dan $a, b \in A$. Maka $(a) + (b) = (d)$ untuk suatu
$d$, yaitu sebuah \emph{faktor persekutuan terbesar}\index{faktor persekutuan
terbesar}: $d \mid a$, $d \mid b$, dan setiap pembagi persekutuan $a, b$
membagi $d$; lebih lanjut $d = au + bv$ untuk suatu $u, v$
(Bézout). Akibatnya setiap \omterm{def:b3:rings:divisibility}{unsur tak tereduksi} di dalam \omterm{def:b3:rings:pidufd}{DIU} bersifat
prima.
\end{lemma}

\begin{proof}
$(a) + (b)$ adalah \omterm{def:b3:rings:ideal}{ideal}, jadi berbentuk $(d)$; dari $a, b \in (d)$ diperoleh $d \mid
a, b$; dan $d = au + bv \in (a) + (b)$. Pembagi persekutuan $c$ dari
$a, b$ membagi $au + bv = d$.

Lema \omterm{def:b3:rings:pidufd}{Euclid}: misalkan $p$ \omterm{def:b3:rings:divisibility}{tak tereduksi}, $p \mid ab$, $p \nmid a$. Sebuah FPB
$d$ dari $p$ dan $a$ membagi $p$, jadi $d$ berupa unit atau \omterm{def:b3:rings:divisibility}{sekawan}
dengan $p$ (karena ketaktereduksian); kesekawanan tersingkir oleh $p \nmid a$. Jadi
$1 = pu + av$, sehingga $b = pub + abv$, dan $p$ membagi kedua sukunya:
jadi $p \mid b$.
\end{proof}

\begin{theorem}\label{thm:b3:rings:pidufd}
\omterm{def:b3:rings:ideal}{Ideal} utama $\Rightarrow$ faktorial.
\end{theorem}

\begin{proof}
\emph{Keberadaan.} Andaikan ada unsur tak nol dan bukan unit $a$ yang tak punya
faktorisasi atas \omterm{def:b3:rings:divisibility}{unsur tak tereduksi}. Maka $a$ sendiri \omterm{def:b3:rings:divisibility}{tak tereduksi}: $a =
a_1b_1$ dengan kedua faktornya bukan unit; paling sedikit salah satunya, katakanlah
$a_1$, kembali tak punya faktorisasi (hasil kali dua unsur yang terfaktorkan
juga terfaktorkan). Dengan mengiterasikannya kita memperoleh $a = a_0, a_1, a_2,
\dots$, masing-masing pembagi sejati pendahulunya dan tanpa faktorisasi,
jadi $(a_0) \subsetneq (a_1) \subsetneq (a_2) \subsetneq \cdots$
--- inklusinya sejati sebab $a_n = a_{n+1}c$ dengan $c$ bukan
unit, sedangkan $(a_n) = (a_{n+1})$ akan memaksa $c \in A^\times$
(coret di daerah integral). Ini bertentangan dengan \cref{lem:b3:rings:acc}.

\emph{Ketunggalan.} Misalkan $p_1 \cdots p_r = q_1 \cdots q_s$ dengan semua
faktornya \omterm{def:b3:rings:divisibility}{tak tereduksi} dan $r \leq s$; kita berinduksi pada $r$. \omterm{def:b3:rings:divisibility}{Unsur prima}
(\cref{lem:b3:rings:bezout}) $p_1$ membagi ruas kanan, sehingga
membagi suatu $q_j$; nomori ulang sehingga $j = 1$. Karena $q_1$ \omterm{def:b3:rings:divisibility}{tak tereduksi} dan
$p_1$ bukan unit, $q_1 = u p_1$ dengan $u \in A^\times$: jadi $p_1, q_1$
\omterm{def:b3:rings:divisibility}{sekawan}. Coretlah $p_1$: $p_2 \cdots p_r = (u q_2)
q_3\cdots q_s$ lalu simpulkan dengan induksi ($r = 1$ memaksa $s = 1$:
sebab unit dikalikan \omterm{def:b3:rings:divisibility}{unsur tak tereduksi} tak mungkin sama dengan $1$).
\end{proof}

\begin{remark}\label{rem:b3:rings:ufdgcd}
Di dalam \omterm{def:b3:rings:pidufd}{DFT}, FPB selalu ada (ambil eksponen terkecil pada
faktorisasinya) dan lema \omterm{def:b3:rings:pidufd}{Euclid} berlaku --- \omterm{def:b3:rings:divisibility}{tak tereduksi} $=$ prima
(\cref{exo:b3:rings:2}) --- tetapi Bézout dapat gagal: di $\Z[X]$,
$\gcd(2, X) = 1$ namun $1 \neq 2U + XV$ (masukkan $X = 0$:
$1 = 2U(0)$, mustahil). Identitas Bézout adalah milik khusus
\omterm{def:b3:rings:pidufd}{DIU} semata.
\end{remark}

\begin{example}[Sebuah ring tanpa faktorisasi tunggal]
\label{ex:b3:rings:nonufd}
Tidak satu pun implikasi \omterm{def:b3:rings:pidufd}{Euclid} $\Rightarrow$ \omterm{def:b3:rings:pidufd}{DIU}
$\Rightarrow$ \omterm{def:b3:rings:pidufd}{DFT} merupakan kesetaraan, dan kegagalan pada
implikasi terakhir layak dilihat sekali secara lengkap. Di dalam
\[
A = \Z[\iu\sqrt5] = \{a + \iu b\sqrt5 : a, b \in \Z\},
\qquad N(a + \iu b\sqrt5) = a^2 + 5b^2,
\]
normanya multiplikatif dan $N(z) = 1$ jika dan hanya jika $z \in
A^\times = \{\pm1\}$. Tinjau
\[
6 = 2 \cdot 3 = (1 + \iu\sqrt5)(1 - \iu\sqrt5).
\]
Keempat faktornya \omterm{def:b3:rings:divisibility}{tak tereduksi}: normanya $4, 9, 6,
6$, dan faktorisasi sejati $z = z_1z_2$ akan memaksa
$N(z_1) \in \{2, 3\}$ --- padahal $a^2 + 5b^2$ tak pernah sama dengan $2$
atau $3$ ($b = 0$ menyisakan $2, 3$ yang bukan kuadrat; $\abs b \geq
1$ memberikan $\geq 5$). Namun $2$ tidak \omterm{def:b3:rings:divisibility}{sekawan} dengan $1 \pm
\iu\sqrt5$ (normanya $4 \neq 6$): jadi ada dua faktorisasi $6$ atas
\omterm{def:b3:rings:divisibility}{unsur tak tereduksi} yang benar-benar berbeda. Setara dengan itu,
di sini \omterm{def:b3:rings:divisibility}{tak tereduksi} $\neq$ prima: $2$ membagi hasil kali $(1 +
\iu\sqrt5)(1 - \iu\sqrt5) = 6$ tetapi tidak membagi satu pun faktornya (lihat
normanya lagi). Perbaikan kegagalan ini secara \omterm{def:b3:rings:ideal}{ideal} ---
memfaktorkan \emph{\omterm{def:b3:rings:ideal}{ideal}} alih-alih unsur --- adalah
kelahiran teori bilangan aljabar; pada tingkat kita, contoh ini
mengukur betapa istimewanya ring \omterm{def:b3:rings:pidufd}{Euclid} $\Z$, $K[X]$,
$\Z[\iu]$ pada bab ini.
\end{example}

\begin{method}\label{met:b3:rings:quotient}
Untuk mengenali \omterm{def:b3:rings:ideal}{ring kuosien} $A/I$, carilah morfisma surjektif
$f \colon A \to B$ berkernel $I$ lalu panggil
\cref{thm:b3:rings:firstiso}; bila $A = C[X]$ berupa ring
polinomial, biasanya $f$ berupa pemasukan nilai. Dengan demikian $\Z[X]/(X^2+1) \cong
\Z[\iu]$ (masukkan $\iu$), $K[X,Y]/(Y - X^2) \cong K[X]$
(ganti $Y$ dengan $X^2$), $\R[X]/(X^2+1) \cong \C$. Untuk menunjukkan $I$
prima atau \omterm{def:b3:rings:primemaximal}{maksimal}, tunjukkan bahwa kuosiennya daerah integral atau lapangan
(\cref{prop:b3:rings:primemaximal}).
\end{method}

\section{Polinomial atas sebuah DFT: Gauss dan Eisenstein}

Sepanjang bagian ini $A$ adalah \omterm{def:b3:rings:pidufd}{DFT} dengan lapangan pecahan $K$
(yang dibangun sebagai lapangan hasil bagi formal $a/b$, $b \neq 0$,
persis seperti $\Q$ dari $\Z$; jilid Tahun ke-2 melakukan
konstruksi itu untuk $\Q$, dan ia berpindah kata demi kata). Tujuan kita:
kefaktorialan berpindah dari $A$ ke $A[X]$, dan ketaktereduksian atas
$A$ pada dasarnya sama dengan ketaktereduksian atas lapangan $K$ yang lebih besar.

\begin{definition}\label{def:b3:rings:content}
Yang disebut \emph{kandungan}\index{kandungan polinomial} $c(P)$ dari
$P \in A[X]$ yang tak nol adalah FPB koefisiennya (tertentu sampai pada
unit); $P$ disebut \emph{primitif} jika $c(P) \in A^\times$. Setiap $P
\in A[X]$ ditulis $P = c(P)\,P_1$ dengan $P_1$ primitif, dan setiap
$P \in K[X]\setminus\{0\}$ ditulis $P = \lambda P_1$ dengan $\lambda
\in K^\times$ dan $P_1 \in A[X]$ primitif (samakan penyebutnya,
lalu keluarkan kandungannya).
\end{definition}

\begin{lemma}[Gauss]\label{lem:b3:rings:gauss}
Hasil kali dua polinomial primitif di $A[X]$ juga primitif;
akibatnya $c(PQ) = c(P)c(Q)$ sampai pada unit.\index{lema Gauss}
\end{lemma}

\begin{proof}
Misalkan $P, Q$ primitif dan andaikan suatu \omterm{def:b3:rings:divisibility}{unsur tak tereduksi} (= prima, karena
\omterm{def:b3:rings:pidufd}{DFT}) $p$ membagi semua koefisien $PQ$. Reduksikan modulo $p$: di
$(A/(p))[X]$ berlaku $\bar P \bar Q = 0$. Padahal $A/(p)$ daerah integral
(sebab $(p)$ prima), sehingga $(A/(p))[X]$ juga daerah integral (koefisien utamanya
saling mengalikan), dan itu memaksa $\bar P = 0$ atau $\bar Q = 0$: berarti $p$ membagi semua
koefisien $P$ atau semua koefisien $Q$, bertentangan dengan keprimitifan. Untuk
akibatnya, tulis $P = c(P)P_1$, $Q = c(Q)Q_1$: maka $PQ =
c(P)c(Q) P_1Q_1$ dengan $P_1Q_1$ primitif.
\end{proof}

\begin{theorem}\label{thm:b3:rings:gaussufd}
Misalkan $A$ sebuah \omterm{def:b3:rings:pidufd}{DFT} dengan lapangan pecahan $K$.
\begin{enumerate}
  \item Polinomial primitif $P \in A[X]$ berderajat $\geq 1$ \omterm{def:b3:rings:divisibility}{tak tereduksi}
        di $A[X]$ jika dan hanya jika ia \omterm{def:b3:rings:divisibility}{tak tereduksi} di $K[X]$.
  \item $A[X]$ adalah \omterm{def:b3:rings:pidufd}{DFT}; \omterm{def:b3:rings:divisibility}{unsur tak tereduksinya} adalah \omterm{def:b3:rings:divisibility}{unsur tak tereduksi}
        $A$ dan polinomial primitif yang \omterm{def:b3:rings:divisibility}{tak tereduksi} atas $K$.
        Khususnya $\Z[X]$, dan lewat induksi $K[X_1, \dots, X_n]$
        serta $\Z[X_1, \dots, X_n]$, semuanya \omterm{def:b3:rings:pidufd}{DFT}.
\end{enumerate}
\end{theorem}

\begin{proof}
(1) ($\Leftarrow$) Jika $P = QR$ di $A[X]$ dengan $Q, R$ bukan unit,
maka tak satu pun faktornya konstan (faktor konstan dari polinomial
primitif pastilah unit), sehingga faktorisasinya sejati di $K[X]$.
($\Rightarrow$) Andaikan $P = QR$ dengan $Q, R \in K[X]$ berderajat
$\geq 1$. Tulis $Q = \lambda Q_1$, $R = \mu R_1$ dengan $Q_1, R_1
\in A[X]$ primitif: maka $P = \lambda\mu\, Q_1R_1$, dan $Q_1R_1$
primitif menurut lema Gauss. Dengan mengambil \omterm{def:b3:rings:content}{kandungannya}, $\lambda\mu \in A^\times$
(kedua ruas berkandungan unit; secara formal $\lambda\mu = c(P) \in
A^\times$ sampai pada unit, dan khususnya $\lambda \mu \in A$): jadi $P
= (\lambda\mu Q_1) R_1$ merupakan faktorisasi sejati di $A[X]$.

(2) Keberadaan: diberikan $P \neq 0$ yang bukan unit, faktorkan $P = c(P)P_1$,
faktorkan $c(P)$ atas \omterm{def:b3:rings:divisibility}{unsur tak tereduksi} $A$, lalu faktorkan $P_1$ di dalam
\omterm{def:b3:rings:pidufd}{DFT} $K[X]$ menjadi $\prod Q_i$ dengan $Q_i \in K[X]$ \omterm{def:b3:rings:divisibility}{tak tereduksi}; dengan menulis
$Q_i = \lambda_i R_i$ dengan $R_i \in A[X]$ primitif (sehingga
\omterm{def:b3:rings:divisibility}{tak tereduksi} atas $K$, sehingga juga di $A[X]$ menurut (1)), hasil kali
$\prod \lambda_i$ merupakan unit $A$ seperti sebelumnya, dan $P_1 = u\prod
R_i$. Ketunggalan: bandingkan bagian konstan dan bagian polinomial sebuah
faktorisasi; konstantanya berkalian menjadi $c(P)$ (lema Gauss), tunggal
menurut kefaktorialan $A$; sedangkan bagian polinomialnya memberikan dua
faktorisasi di $K[X]$ atas polinomial yang sama, sehingga keduanya bersesuaian sampai pada
konstanta $K^\times$ (kefaktorialan $K[X]$,
\cref{thm:b3:rings:pidufd}), dan polinomial primitif yang bersesuaian dan
\omterm{def:b3:rings:divisibility}{sekawan} di $K[X]$ juga \omterm{def:b3:rings:divisibility}{sekawan} di $A[X]$: sebab jika $R = \lambda R'$
dengan $R, R'$ primitif dan $\lambda \in K^\times$, maka mengambil
\omterm{def:b3:rings:content}{kandungannya} memaksa $\lambda \in A^\times$.
\end{proof}

\begin{theorem}[Kriteria ketaktereduksian]\label{thm:b3:rings:criteria}
Misalkan $A$ sebuah \omterm{def:b3:rings:pidufd}{DFT}, $K$ lapangan pecahannya, dan $P = a_nX^n + \dots
+ a_0 \in A[X]$ primitif berderajat $n \geq 1$.
\begin{enumerate}
  \item (\emph{Reduksi}) Jika $p \in A$ prima, $p \nmid a_n$,
        dan reduksi $\bar P$ \omterm{def:b3:rings:divisibility}{tak tereduksi} di
        $(A/(p))[X]$, maka $P$ \omterm{def:b3:rings:divisibility}{tak tereduksi} di $K[X]$ (sehingga juga di
        $A[X]$).
  \item (\emph{Eisenstein}\index{kriteria Eisenstein}) Jika ada
        bilangan prima $p$ dengan $p \nmid a_n$, $p \mid a_i$ untuk $0
        \leq i < n$, dan $p^2 \nmid a_0$, maka $P$ \omterm{def:b3:rings:divisibility}{tak tereduksi}
        di $K[X]$ (sehingga juga di $A[X]$).
\end{enumerate}
\end{theorem}

\begin{proof}
Menurut \cref{thm:b3:rings:gaussufd}(1), faktorisasi sejati atas
$K$ melahirkan $P = QR$ dengan $Q, R \in A[X]$, $\deg Q, \deg R \geq 1$
(konstanta tersingkir: ia akan menjadi unit atau merusak
keprimitifan).

(1) Reduksikan modulo $p$: $\bar P = \bar Q\bar R$ di $(A/(p))[X]$. Karena
$p \nmid a_n$ dan $\deg$ hanya dapat turun di bawah reduksi, $\deg \bar
Q = \deg Q \geq 1$ dan $\deg\bar R = \deg R \geq 1$ (koefisien utamanya
berkalian menjadi $\bar a_n \neq 0$, jadi tak satu pun turun):
jadi $\bar P$ terfaktorkan secara sejati --- kontradiksi.

(2) Reduksikan modulo $p$: $\bar Q \bar R = \bar P = \bar a_n X^n$ (semua
koefisien yang lebih rendah lenyap). Di daerah integral $(A/(p))[X]$,
faktorisasi $cX^n$ ($c \ne 0$) hanya berupa konstanta dan pangkat
murni $c'X^k$: memang jika $\bar Q\bar R = \bar a_nX^n$ dan misalnya
$\bar Q$ mempunyai koefisien tak nol berderajat $< \deg\bar Q$, ambil
suku tak nol terendah: $\operatorname{val}(\bar Q\bar R) =
\operatorname{val}\bar Q + \operatorname{val}\bar R$ (di daerah integral),
yang harus sama dengan $n = \deg\bar Q + \deg\bar R$, sehingga memaksa
$\operatorname{val} = \deg$ pada keduanya: jadi keduanya monomial. Seperti
di atas, derajatnya tidak turun, sehingga suku konstan $Q$ dan $R$,
yaitu $Q(0), R(0)$, habis dibagi $p$ --- keduanya, sebab kedua
reduksinya monomial berderajat $\geq 1$. Maka $p^2 \mid
Q(0)R(0) = a_0$: kontradiksi.
\end{proof}

\begin{example}\label{ex:b3:rings:eisenstein}
$X^n - p$ \omterm{def:b3:rings:divisibility}{tak tereduksi} atas $\Q$ untuk setiap bilangan prima $p$ dan $n
\geq 1$ (\omterm{thm:b3:rings:criteria}{Eisenstein} pada $p$): jadi ada polinomial \omterm{def:b3:rings:divisibility}{tak tereduksi}
berderajat berapa pun atas $\Q$ --- sangat berbeda dengan $\C$ (hanya derajat
$1$, menurut d'Alembert--Gauss, yang dibuktikan pada \cref{ch:b3:holomorphic}) dan
$\R$ (derajat $1, 2$). Muslihat \emph{penggeseran} memperluas
jangkauan \omterm{thm:b3:rings:criteria}{Eisenstein}: \emph{polinomial siklotomik}\index{polinomial
siklotomik} ke-$p$, yaitu $\Phi_p = X^{p-1} +
\dots + X + 1 = \frac{X^p - 1}{X - 1}$, mempunyai
\[
\Phi_p(X + 1) = \frac{(X+1)^p - 1}{X}
= X^{p-1} + \binom{p}{1}X^{p-2} + \dots + \binom{p}{p-1},
\]
yang memenuhi \omterm{thm:b3:rings:criteria}{Eisenstein} pada $p$ (sebab $p \mid \binom pk$ untuk $0 < k < p$, dan
$\binom{p}{p-1} = p \not\equiv 0 \bmod p^2$): jadi $\Phi_p(X+1)$, dan karenanya
$\Phi_p$, \omterm{def:b3:rings:divisibility}{tak tereduksi} atas $\Q$. Inilah jantung aljabar dari
kisah segi-$17$ yang dituturkan pada \cref{ch:b3:galois}.
\end{example}

\begin{method}\label{met:b3:rings:irreducibility}
Untuk membuktikan $P \in \Z[X]$ \omterm{def:b3:rings:divisibility}{tak tereduksi} atas $\Q$: (i) jadikan $P$
primitif; (ii) cobalah \omterm{thm:b3:rings:criteria}{Eisenstein}, pada $P(X)$ dan pada geseran $P(X \pm
1)$; (iii) cobalah reduksi modulo bilangan prima kecil yang tak membagi
koefisien utamanya --- ketaktereduksian modulo \emph{satu} $p$
sudah cukup, dan atas $\mathbb F_p$ ketaktereduksian merupakan pemeriksaan berhingga
(tanpa akar berarti tanpa faktor berderajat $1$; lalu ujilah faktor yang berhingga
banyaknya pada setiap derajat $\leq \deg P/2$); (iv) jika semuanya gagal,
pakailah koefisien tak tentu. Awas: keterreduksian modulo setiap $p$
\emph{tidak} mengakibatkan keterreduksian atas $\Q$
(\cref{exo:b3:rings:11}).
\end{method}

\section{Ring Noether}

\begin{definition}\label{def:b3:rings:noetherian}
Sebuah ring $A$ disebut \emph{Noether}\index{ring Noether} jika setiap
\omterm{def:b3:rings:ideal}{ideal} $A$ dibangun secara berhingga.
\end{definition}

\begin{proposition}\label{prop:b3:rings:noethacc}
$A$ bersifat \omterm{def:b3:rings:noetherian}{Noether} jika dan hanya jika setiap barisan naik \omterm{def:b3:rings:ideal}{ideal}
akhirnya konstan (\emph{syarat rantai naik}), dan jika dan hanya jika setiap
keluarga \omterm{def:b3:rings:ideal}{ideal} yang tak kosong mempunyai unsur \omterm{def:b3:rings:primemaximal}{maksimal} (menurut inklusi).
\end{proposition}

\begin{proof}
\emph{(Dibangun berhingga $\Rightarrow$ syarat rantai naik)}: untuk rantai $I_1 \subseteq I_2
\subseteq \cdots$, gabungannya $I$ adalah \omterm{def:b3:rings:ideal}{ideal} yang dibangun oleh $x_1,
\dots, x_r$; semua $x_i$ terletak pada suatu $I_N$, jadi $I = I_N = I_n$ untuk
$n \geq N$. \emph{(Syarat rantai naik $\Rightarrow$ unsur \omterm{def:b3:rings:primemaximal}{maksimal})}: seandainya
keluarga tak kosong $\mathcal F$ tak punya unsur \omterm{def:b3:rings:primemaximal}{maksimal}, pilihlah $I_1 \in
\mathcal F$, lalu secara induktif $I_{n+1} \supsetneq I_n$ di
$\mathcal F$ (mungkin dilakukan karena $I_n$ tidak \omterm{def:b3:rings:primemaximal}{maksimal}): terbentuklah rantai naik
sejati yang tak berujung. (Ini memakai aksioma pilihan bergantung,
bentuk lemah dari aksioma pilihan yang tak kita persoalkan.)
\emph{(Unsur \omterm{def:b3:rings:primemaximal}{maksimal} $\Rightarrow$ dibangun berhingga)}: diberikan sebuah \omterm{def:b3:rings:ideal}{ideal} $I$,
keluarga \omterm{def:b3:rings:ideal}{ideal} yang dibangun berhingga dan termuat di $I$ bersifat
tak kosong (ada $(0)$); unsur \omterm{def:b3:rings:primemaximal}{maksimalnya} $J = (x_1, \dots, x_r)$ pastilah
sama dengan $I$: sebab kalau tidak, menambahkan $x \in I \setminus J$ pada
pembangunnya akan menghasilkan anggota keluarga yang lebih besar secara sejati.
\end{proof}

\begin{theorem}[Teorema basis Hilbert]\label{thm:b3:rings:hilbert}
Jika $A$ bersifat \omterm{def:b3:rings:noetherian}{Noether}, maka $A[X]$ pun demikian. Karena itu $A[X_1, \dots,
X_n]$ dan setiap kuosiennya juga demikian.\index{teorema basis Hilbert}
\end{theorem}

\begin{proof}
Misalkan $I$ sebuah \omterm{def:b3:rings:ideal}{ideal} $A[X]$, dan andaikan $I$ tidak dibangun secara
berhingga. Bangunlah sebuah barisan: $f_1 \in I \setminus \{0\}$ yang
berderajat minimal, lalu secara induktif $f_{k+1} \in I \setminus (f_1,
\dots, f_k)$ yang berderajat minimal (himpunannya tak kosong menurut
pengandaian). Derajat $d_k = \deg f_k$ tidak pernah turun (menurut
keminimalan setiap pilihan: $f_{k+1}$ sudah tersedia pada langkah
$k+1$\dots\ lebih tepatnya, $f_{k+1} \notin (f_1,\dots,f_k) \supseteq
(f_1, \dots, f_{k-1})$, jadi $f_{k+1}$ ikut bersaing pada langkah $k$ dan kalah
atau seri: $d_{k+1} \geq d_k$). Misalkan $a_k \in A$ koefisien utama
$f_k$. Rantai \omterm{def:b3:rings:ideal}{ideal} $(a_1) \subseteq (a_1,
a_2) \subseteq \cdots$ stabil: $a_{n+1} \in (a_1, \dots,
a_n)$ untuk suatu $n$, katakanlah $a_{n+1} = \sum_{k\leq n} u_k a_k$.
Tinjau
\[
g = f_{n+1} - \sum_{k=1}^{n} u_k X^{\,d_{n+1} - d_k} f_k .
\]
Maka $g \in I \setminus (f_1, \dots, f_n)$ (jumlahnya terletak di
\omterm{def:b3:rings:ideal}{ideal} itu, sedangkan $f_{n+1}$ tidak), namun koefisien berderajat
$d_{n+1}$ saling meniadakan: $\deg g < d_{n+1}$, dan itu bertentangan dengan
keminimalan $d_{n+1} = \deg f_{n+1}$.

Dengan mengiterasikannya, $A[X_1, \dots, X_n] = (A[X_1, \dots, X_{n-1}])[X_n]$
bersifat \omterm{def:b3:rings:noetherian}{Noether}; sedangkan kuosien $A/I$ bersifat \omterm{def:b3:rings:noetherian}{Noether} karena \omterm{def:b3:rings:ideal}{idealnya}
$J/I$ terangkat menjadi \omterm{def:b3:rings:ideal}{ideal} $A$ (korespondensi), yang pembangunnya
berhingga banyaknya dan terproyeksikan menjadi pembangun.
\end{proof}

\begin{remark}\label{rem:b3:rings:noethmeaning}
Sifat \omterm{def:b3:rings:noetherian}{Noether} adalah aksioma keberhinggaan pada geometri aljabar: setiap
sistem persamaan polinomial dengan $n$ variabel, betapa pun tak berhingga,
setara dengan sistem yang berhingga banyaknya --- himpunan penyelesaiannya terpotong
oleh polinomial yang berhingga banyaknya. Setiap \omterm{def:b3:rings:pidufd}{DIU} bersifat \omterm{def:b3:rings:noetherian}{Noether}
(secara trivial); sedangkan $\Z[X_1, X_2, \dots]$ dengan variabel tak berhingga banyaknya
tidak ($(X_1) \subsetneq (X_1, X_2) \subsetneq \cdots$).
Ring yang bukan \omterm{def:b3:rings:noetherian}{Noether} juga muncul secara alami dalam analisis: fungsi
kontinu pada $\intcc01$ membentuk salah satunya (\cref{exo:b3:rings:10}).
\end{remark}

\section{Latihan}

\begin{exercise}[$\star$]\label{exo:b3:rings:1}
Kenali kuosien berikut: (a) $\Z[X]/(X^2 + 1) \cong \Z[\iu]$;
(b) $\R[X]/(X^2+1) \cong \C$; (c) $\mathbb F_2[X]/(X^2 + X + 1)$
adalah lapangan dengan $4$ unsur --- tuliskan tabel perkaliannya.
\end{exercise}

\begin{exercise}[$\star$]\label{exo:b3:rings:2}
(a) Tunjukkan bahwa di dalam \omterm{def:b3:rings:pidufd}{DFT} setiap \omterm{def:b3:rings:divisibility}{unsur tak tereduksi} bersifat prima.
(b) Tunjukkan bahwa daerah integral yang berhingga adalah lapangan.
(c) Simpulkan bahwa di dalam ring berhingga setiap \omterm{def:b3:rings:primemaximal}{ideal prima} bersifat \omterm{def:b3:rings:primemaximal}{maksimal}.
\end{exercise}

\begin{exercise}[$\star$]\label{exo:b3:rings:3}
Di $\Z[\iu\sqrt5]$: periksalah bahwa $3$, $1 + \iu\sqrt5$ dan $1 -
\iu\sqrt5$ \omterm{def:b3:rings:divisibility}{tak tereduksi}, bahwa $9 = 3\cdot 3 = (2 +
\iu\sqrt5)(2 - \iu\sqrt5)$, lalu simpulkan sekali lagi (setelah
\cref{prop:b3:rings:primeirred}) bahwa $\Z[\iu\sqrt5]$ bukan
\omterm{def:b3:rings:pidufd}{DFT}. Di manakah tepatnya ketunggalan itu gagal?
\end{exercise}

\begin{exercise}[$\star\star$]\label{exo:b3:rings:4}
(a) Tunjukkan bahwa $\Z[\iu]$ bersifat \omterm{def:b3:rings:pidufd}{Euclid} terhadap norma $N(x + \iu y) =
x^2 + y^2$: diberikan $a, b \neq 0$, pilihlah $q \in \Z[\iu]$ yang terdekat
dengan $a/b \in \C$.
(b) Tentukan $\Z[\iu]^\times$.
(c) Pertanyaan yang sama untuk $\Z[\iu\sqrt2]$ dan $N(x + \iu y\sqrt2) =
x^2 + 2y^2$. Mengapa argumen yang sama gagal untuk
$\Z[\iu\sqrt5]$?
\end{exercise}

\begin{exercise}[$\star\star$]\label{exo:b3:rings:5}
Misalkan $A$ sebuah ring. (a) Tunjukkan bahwa jika $x$ nilpoten ($x^n = 0$
untuk suatu $n$) maka $1 + x \in A^\times$. (b) Tunjukkan bahwa jika $A$
daerah integral, maka $A[X]^\times = A^\times$; berilah contoh penyangkal atas
$\Z/4\Z$. (c) Tunjukkan bahwa daerah integral tak punya unsur idempoten ($e^2 = e$)
selain $0, 1$, dan tak punya unsur nilpoten selain $0$.
\end{exercise}

\begin{exercise}[$\star\star$]\label{exo:b3:rings:6}
Di $A = K[X, Y]$: (a) tunjukkan bahwa \omterm{def:b3:rings:ideal}{ideal} $(X, Y)$ \omterm{def:b3:rings:primemaximal}{maksimal} tetapi
bukan \omterm{def:b3:rings:ideal}{ideal} utama --- sehingga $K[X,Y]$ merupakan \omterm{def:b3:rings:pidufd}{DFT}
(\cref{thm:b3:rings:gaussufd}) yang bukan \omterm{def:b3:rings:pidufd}{DIU}; (b) kenali
$K[X, Y]/(Y - X^2)$ dan $K[X,Y]/(XY - 1)$ sebagai subring fungsi
rasional; (c) apakah $(Y - X^2)$ prima? \omterm{def:b3:rings:primemaximal}{maksimal}?
\end{exercise}

\begin{exercise}[$\star\star$]\label{exo:b3:rings:7}
\omterm{def:b3:rings:divisibility}{Tak tereduksi} atau tidak atas $\Q$: $X^5 - 12X^3 + 36X - 12$;\quad
$X^4 + X + 1$ \emph{(reduksikan modulo $2$)};\quad $X^4 + 4$;\quad
$\Phi_8 = X^4 + 1$ \emph{(geser sejauh $1$)};\quad $X^3 - X - 1$.
\end{exercise}

\begin{exercise}[$\star\star$]\label{exo:b3:rings:8}
(a) Dari \cref{thm:b3:rings:crt}, buktikan bahwa fungsi Euler bersifat
multiplikatif pada argumen yang relatif prima dan bahwa $\varphi(p^k) =
p^{k-1}(p-1)$; peroleh kembali $\varphi(n) = n\prod_{p \mid n}(1 -
\frac1p)$.
(b) Selesaikan: $x \equiv 2 \pmod 7$, $x \equiv 5 \pmod{11}$, $x
\equiv 1 \pmod{13}$, sambil menampilkan unsur idempoten $e_k$ pada
bukti \cref{thm:b3:rings:crt}.
\end{exercise}

\begin{exercise}[$\star\star\star$]\label{exo:b3:rings:9}
(Nilradikal) Misalkan $\operatorname{Nil}(A)$ himpunan semua
unsur nilpoten. (a) Tunjukkan bahwa $\operatorname{Nil}(A)$ adalah
\omterm{def:b3:rings:ideal}{ideal} yang termuat di setiap \omterm{def:b3:rings:primemaximal}{ideal prima}. (b) Sebaliknya, misalkan $a$
bukan nilpoten; dengan memakai lema Zorn pada \omterm{def:b3:rings:ideal}{ideal} yang menghindari $S =
\{a^n : n \in \N\}$, hasilkanlah \omterm{def:b3:rings:primemaximal}{ideal prima} yang tak memuat $a$.
Simpulkan:
\[
\operatorname{Nil}(A) = \bigcap_{\mathfrak p \text{ prima}}
\mathfrak p .
\]
\end{exercise}

\begin{exercise}[$\star\star\star$]\label{exo:b3:rings:10}
(a) Misalkan $A$ bersifat \omterm{def:b3:rings:noetherian}{Noether} dan $f \colon A \to A$ morfisma ring
yang surjektif. Tunjukkan bahwa $f$ injektif. \emph{(Tinjaulah $\ker
f \subseteq \ker f^2 \subseteq \cdots$.)}
(b) Tunjukkan bahwa ring $\mathcal C(\intcc01, \R)$ yang beranggotakan fungsi
kontinu tidak bersifat \omterm{def:b3:rings:noetherian}{Noether}. \emph{(Tinjaulah $I_n = \{f : f = 0
\text{ pada } \intcc0{1/n}\}$.)}
(c) Tunjukkan bahwa di dalam \emph{daerah integral} yang \omterm{def:b3:rings:noetherian}{Noether}, setiap unsur tak nol
yang bukan unit merupakan hasil kali berhingga \omterm{def:b3:rings:divisibility}{unsur tak tereduksi} --- sehingga
ketakfaktorialan $\Z[\iu\sqrt 5]$ semata-mata kegagalan
\emph{ketunggalan}.
\end{exercise}

\begin{exercise}[$\star\star\star$]\label{exo:b3:rings:11}
Misalkan $P = X^4 + 1$.
(a) Tunjukkan bahwa $P$ \omterm{def:b3:rings:divisibility}{tak tereduksi} atas $\Q$
(\cref{exo:b3:rings:7}).
(b) Tunjukkan bahwa $P$ terreduksi modulo \emph{setiap} bilangan prima $p$:
tangani $p = 2$; lalu, untuk $p$ ganjil, tunjukkan bahwa $8 \mid p^2 - 1$ dan
terimalah untuk sementara (dibuktikan pada \cref{ch:b3:galois}) bahwa
grup multiplikatif dari lapangan dengan $p^2$ unsur bersifat siklik,
lalu simpulkan bahwa $P$ terpecah atas dua faktor kuadrat modulo $p$;
tuliskan secara eksplisit bila salah satu dari $-1$, $2$, $-2$ merupakan kuadrat modulo
$p$, dan tunjukkan bahwa selalu ada salah satunya.
\end{exercise}

\begin{exercise}[$\star\star$]\label{exo:b3:rings:12}
(Unsur idempoten memecah ring) Sebuah unsur $e$ dari ring komutatif
$A$ disebut \emph{idempoten} jika $e^2 = e$.
(a) Tunjukkan bahwa jika $e$ idempoten maka $1 - e$ pun idempoten, dan bahwa
pemetaan $x \mapsto (ex, (1-e)x)$ adalah isomorfisma ring $A
\cong Ae \times A(1-e)$, dengan $Ae$ berupa ring bersatuan $e$.
(b) Carilah semua unsur idempoten dari daerah integral, dan dari $\Z/12\Z$;
tampilkan isomorfisma $\Z/12\Z \cong \Z/4\Z \times \Z/3\Z$
dengan menyebutkan kedua unsur idempotennya yang tak trivial.
(c) Tunjukkan bahwa penguraian $\Z/n\Z$ lewat teorema sisa Tiongkok
(\cref{ex:b3:rings:crtz}) bersesuaian persis dengan
unsur idempoten $e_i \equiv 1 \bmod p_i^{a_i}$, $e_i \equiv 0$
modulo pangkat prima lainnya: ring terurai menurut unsur
idempotennya sebagaimana ruang terurai menurut proyeksinya.
\end{exercise}

\section{Soal: teorema dua kuadrat Fermat}

\begin{problem}[{Soal akhir pekan --- jumlah dua kuadrat, lewat
$\Z[\iu]$}]\label{pb:b3:rings:1}
Bilangan bulat manakah yang merupakan jumlah dua kuadrat? Jawaban Fermat (1640) adalah
salah satu permata aritmetika; bilangan bulat Gauss mengubah pembuktiannya
menjadi teori ring. Sepanjang soal ini, $N(x + \iu y) = x^2 + y^2$ melambangkan
normanya, $\Z[\iu]$ bersifat \omterm{def:b3:rings:pidufd}{Euclid} (\cref{exo:b3:rings:4}) sehingga
menjadi \omterm{def:b3:rings:pidufd}{DIU} sekaligus \omterm{def:b3:rings:pidufd}{DFT}, dan \emph{prima Gauss} berarti \omterm{def:b3:rings:divisibility}{unsur prima} (=
\omterm{def:b3:rings:divisibility}{tak tereduksi}) di $\Z[\iu]$.

\medskip
\noindent\textbf{Bagian I --- Norma dan prima Gauss.}
\begin{enumerate}
  \item Periksalah $N(zw) = N(z)N(w)$, simpulkan kembali $\Z[\iu]^\times =
        \{\pm1, \pm\iu\}$, lalu buktikan \emph{identitas
        Brahmagupta}: hasil kali dua jumlah dua kuadrat juga merupakan jumlah
        dua kuadrat.
  \item Tunjukkan bahwa jika $N(z)$ berupa bilangan prima, maka $z$ merupakan
        prima Gauss.
  \item Tunjukkan bahwa setiap prima Gauss $\pi$ membagi tepat satu
        bilangan prima $p$ \emph{(tinjaulah $N(\pi) = \pi\bar\pi$)},
        dan bahwa lalu $N(\pi) \in \{p, p^2\}$.
  \item Simpulkan dikotominya: untuk setiap bilangan prima $p$, entah $p$
        tetap prima di $\Z[\iu]$ (dan tak ada prima Gauss bernorma
        $p$), entah $p = \pi\bar\pi$ dengan $\pi$ sebuah prima
        Gauss bernorma $p$ --- dan lalu $p = a^2 + b^2$.
\end{enumerate}

\medskip
\noindent\textbf{Bagian II --- Teorema Wilson dan $-1$ modulo
$p$.}
\begin{enumerate}[resume]
  \item Buktikan \emph{teorema Wilson}: untuk $p$ prima, $(p-1)!
        \equiv -1 \pmod p$. \emph{(Pasangkan setiap sisa dengan
        inversnya; mana saja yang berpasangan dengan dirinya sendiri?)}
  \item Misalkan $p$ bilangan prima ganjil dan $m = \frac{p-1}2$. Tunjukkan bahwa
        $(m!)^2 \equiv (-1)^{m+1} \pmod p$ \emph{(di dalam $(p-1)!$,
        gantilah setiap faktor $k > m$ dengan $-(p - k)$)}.
  \item Simpulkan: $-1$ merupakan kuadrat modulo $p$ jika dan hanya jika $p = 2$ atau $p
        \equiv 1 \pmod 4$. \emph{(Untuk arah ``hanya jika'': jika $x^2
        \equiv -1$, berapa orde $x$ di $(\Z/p\Z)^\times$,
        dan apa kata teorema Lagrange?)}
\end{enumerate}

\medskip
\noindent\textbf{Bagian III --- Hukum pemecahan.}
\begin{enumerate}[resume]
  \item Misalkan $p \equiv 1 \pmod 4$, dan $x$ memenuhi $p \mid x^2 + 1 =
        (x + \iu)(x - \iu)$. Tunjukkan bahwa $p$ \emph{bukan}
        prima Gauss, lalu simpulkan bersama Bagian I: $p = a^2 + b^2$.
  \item Misalkan $p \equiv 3 \pmod 4$. Tunjukkan secara langsung bahwa $p$ bukan
        jumlah dua kuadrat \emph{(lihat kuadrat modulo $4$)}, lalu simpulkan
        bahwa $p$ tetap menjadi prima Gauss.
  \item Tuntaskan kasus $p = 2$: tampilkan faktorisasi $2 =
        -\iu(1+\iu)^2$ dan periksa bahwa $1 + \iu$ merupakan prima
        Gauss. (Bilangan $2$ adalah satu-satunya prima yang \emph{bercabang}: habis dibagi
        kuadrat sebuah prima Gauss sampai pada unit.)
  \item Susunlah klasifikasi prima Gauss, sampai pada
        unit: $1 + \iu$; bilangan bulat $p \equiv 3 \pmod 4$; dan
        pasangan \omterm{def:b3:rings:divisibility}{sekawan} $\pi, \bar\pi$ bernorma $p \equiv 1 \pmod
        4$. Periksalah pada $5 = (2+\iu)(2-\iu)$ dan pada $3$.
\end{enumerate}

\medskip
\noindent\textbf{Bagian IV --- Teorema dua kuadrat.}
\begin{enumerate}[resume]
  \item Buktikan separuh yang langsung: jika pada faktorisasi $n =
        \prod p_i^{\alpha_i}$ setiap bilangan prima $\equiv 3 \pmod 4$
        muncul dengan eksponen genap, maka $n$ merupakan jumlah dua
        kuadrat. \emph{(Brahmagupta + Bagian II--III.)}
  \item Buktikan konversnya: jika $n = a^2 + b^2 = N(a + \iu b)$ dan
        $q \equiv 3 \pmod 4$ membagi $n$, tunjukkan bahwa $q$, yang merupakan
        prima Gauss, membagi $a + \iu b$ atau $a - \iu b$, bahwa
        ia sesungguhnya membagi baik $a$ maupun $b$, lalu simpulkan dengan
        induksi pada $n$ bahwa eksponen $q$ di dalam $n$ bernilai genap.
  \item Nyatakan teorema akhirnya. Manakah di antara $2025$, $2026$, $2027$
        yang merupakan jumlah dua kuadrat? \emph{($2025 = 81 \cdot 25$;
        $2026 = 2 \cdot 1013$ dengan $1013$ prima; $2027$ prima.)}
  \item (Penutup) Tunjukkan bahwa bilangan prima $p \equiv 1 \pmod 4$ merupakan
        jumlah dua kuadrat secara pada dasarnya tunggal: jika $p =
        a^2 + b^2 = c^2 + d^2$ (dengan bilangan bulat positif), maka $\{a, b\}
        = \{c, d\}$. \emph{(Ketunggalan faktorisasi di
        $\Z[\iu]$.)}
\end{enumerate}

\medskip
\noindent\textbf{Bagian V --- Mencacah representasi: rumus Jacobi
dan deret Leibniz.} Tulis $r_2(n) = \#\{(a, b) \in
\Z^2 : a^2 + b^2 = n\}$ (pasangan terurut, dengan tanda dan nol ikut
dihitung), dan misalkan $\chi$ karakter tak trivial modulo $4$:
$\chi(d) = +1$ jika $d \equiv 1$, $-1$ jika $d \equiv 3 \pmod4$,
dan $0$ jika $d$ genap.
\begin{enumerate}[resume]
  \item (Pemanasan, sebagai pembanding) Bilangan bulat manakah yang merupakan
        \emph{selisih} dua kuadrat? Tunjukkan: $n = a^2 - b^2$
        dengan $a, b \in \Z$ jika dan hanya jika $n \not\equiv 2 \pmod 4$ ---
        tanpa teori ring sama sekali, dan tanpa struktur yang sebanding dengan
        yang menyusul.
  \item Tunjukkan bahwa $r_2(n)$ adalah banyaknya $z \in \Z[\iu]$ dengan
        $N(z) = n$. Dengan menulis $n = 2^{a}\prod_jp_j^{b_j}
        \prod_kq_k^{c_k}$ dengan $p_j \equiv 1$, $q_k \equiv 3
        \pmod4$, pakailah klasifikasi pada pertanyaan 11 dan
        ketunggalan faktorisasi untuk menunjukkan: $z$ semacam itu ada jika dan hanya jika semua
        $c_k$ genap, dan dalam hal itu
        \[
        r_2(n) = 4\prod_j\,(b_j + 1) .
        \]
        \emph{(Cacahlah: $z = u\,(1+\iu)^{a}\prod_j\pi_j^{s_j}
        \bar\pi_j^{\,b_j - s_j}\prod_kq_k^{c_k/2}$ dengan $u$ sebuah
        unit dan $0 \leq s_j \leq b_j$; mengapa daftar ini
        menyeluruh dan tanpa pengulangan?)}
  \item Tunjukkan bahwa $d \mapsto \chi(d)$ multiplikatif
        sepenuhnya, simpulkan bahwa $n \mapsto
        \sum_{d \mid n}\chi(d)$ multiplikatif, lalu hitunglah nilainya
        pada pangkat prima: nilainya $1$ pada $2^a$; $b + 1$ pada
        $p^b$ ($p \equiv 1$); serta $1$ atau $0$ pada $q^c$ ($q \equiv 3$)
        menurut genap atau ganjilnya $c$.
  \item Simpulkan \emph{teorema Jacobi}:
        \[
        r_2(n) = 4\sum_{d \mid n}\chi(d) =
        4\bigl(d_1(n) - d_3(n)\bigr),
        \]
        dengan $d_i(n)$ mencacah pembagi $\equiv i \pmod 4$.
        Periksalah pada $n = 3, 5, 9, 25$, lalu daftarkan $16$
        representasi dari $65$.
  \item (Lingkaran) Tunjukkan bahwa $\sum_{n \leq x}r_2(n)$ adalah
        banyaknya titik kisi $\Z^2$ di dalam cakram tertutup
        berjari-jari $\sqrt x$, lalu buktikan
        \[
        \sum_{n\leq x}r_2(n) = \pi x + O(\sqrt x)
        \]
        \emph{(setiap titik kisi memiliki satu persegi satuan; bandingkan
        luasnya, dengan galat yang hidup di dalam cincin selebar
        $O(1)$)}.
  \item (Leibniz, dibaca secara aritmetika) Gabungkan pertanyaan 19--20:
        \[
        \sum_{d \leq x}\chi(d)\Bigl\lfloor\frac
        xd\Bigr\rfloor = \frac{\pi x}4 + O(\sqrt x),
        \]
        lalu simpulkan --- setelah fungsi lantainya dilepas dengan hati-hati ---
        deret Leibniz
        \[
        1 - \frac13 + \frac15 - \frac17 + \dots = \frac\pi4 .
        \]
        Deret berganti tanda atas kebalikan bilangan ganjil itu \emph{adalah}
        rata-rata kelebihan pembagi $\equiv 1$ atas pembagi
        $\equiv 3$: analisis yang dihitung oleh aritmetika.
  \item (Seberapa langka jumlah dua kuadrat?) Tunjukkan bahwa tak ada
        bilangan bulat $\equiv 3 \pmod 4$ yang merupakan jumlah dua kuadrat (dengan dua
        cara: kuadrat modulo $4$, atau kriteria keparitasan pada
        pertanyaan 17), sehingga sekurang-kurangnya seperempat bilangan bulat
        terlewatkan; lalu tunjukkan bahwa rata-rata $\frac1x\sum_{n\leq
        x}r_2(n) \to \pi$ pada pertanyaan 20 sejalan dengan
        bilangan bulat terwakilkan yang berkepadatan $0$ ---
        tampilkan bilangan bulat dengan representasi yang luar biasa banyak
        (ambil hasil kali banyak bilangan prima $\equiv 1 \bmod 4$) untuk
        menjelaskan bagaimana proporsi yang menuju nol masih dapat membawa
        rata-rata positif. (Landau membuktikan bahwa kepadatan sejatinya meluruh
        seperti $1/\sqrt{\log x}$; itu di luar jangkauan perkakas kita, tetapi
        mekanismenya kini sudah terlihat.)
\end{enumerate}

\medskip
\noindent\textbf{Bagian VI --- Pelengkap: representasi primitif
dan Pythagoras.}
\begin{enumerate}[resume]
  \item Sebutlah representasi $n = a^2 + b^2$ \emph{primitif}
        jika $\gcd(a, b) = 1$. Tunjukkan bahwa $n \geq 1$ mempunyai
        representasi primitif jika dan hanya jika $4 \nmid n$ dan tak ada bilangan prima
        $q \equiv 3 \pmod 4$ yang membagi $n$. \emph{(Untuk
        keperluannya, pakai lagi penurunan pada pertanyaan 13 dan kuadrat
        modulo $4$; untuk kecukupannya, bangunlah $z$ dari $1 + \iu$
        dan $\pi_j$ saja --- tanpa \omterm{def:b3:rings:divisibility}{sekawannya} --- lalu jelaskan
        mengapa faktor prima persekutuan $a$ dan $b$ akan memaksa
        baik $\pi_j$ maupun $\bar\pi_j$, atau $(1+\iu)^2$, masuk ke
        dalam $z$.)}
  \item (Tripel Pythagoras) Misalkan $a^2 + b^2 = c^2$ dengan $a, b,
        c$ positif, $\gcd(a, b) = 1$ dan $b$ genap. Tunjukkan bahwa
        $a + \iu b$ dan $a - \iu b$ relatif prima di $\Z[\iu]$
        \emph{(pembagi prima Gauss persekutuannya akan membagi
        $2a$ dan $2b$, padahal $c$ ganjil)}, simpulkan dari ketunggalan
        faktorisasi bahwa $a + \iu b = u(m + \iu n)^2$ untuk suatu
        unit $u$, lalu simpulkan parameterisasi klasiknya:
        setelah $a$ dan $b$ ditukar bila perlu,
        \[
        a = m^2 - n^2, \qquad b = 2mn, \qquad c = m^2 + n^2,
        \]
        dengan $m > n \geq 1$ relatif prima dan berkeparitasan berlawanan.
        Peroleh kembali $(3, 4, 5)$ dan $(21, 20, 29)$ dari $(m, n) =
        (2, 1)$ dan $(5, 2)$.
  \item (Pemeriksaan numerik) Ambil $x = 25$. Hitunglah
        $r_2(n)$ untuk $1 \leq n \leq 25$ dari rumus Jacobi,
        periksa bahwa nilai tak nolnya muncul tepat pada $n = 1,
        2, 4, 5, 8, 9, 10, 13, 16, 17, 18, 20, 25$, dan bahwa
        \[
        \sum_{n \leq 25} r_2(n) = 80
        = 4\sum_{d \leq 25}\chi(d)
        \Bigl\lfloor\frac{25}d\Bigr\rfloor .
        \]
        Periksalah bahwa cakram tertutup berjari-jari $5$ memuat $81$
        titik kisi, lalu bandingkan dengan $\pi x \approx 78.5$:
        galatnya masih jauh di dalam $O(\sqrt x)$ pada pertanyaan
        20.

\end{enumerate}
\end{problem}
